\documentclass[letterpaper]{article}
\usepackage{amsmath}
\usepackage{tikz}
\usepackage{epigraph}
\usepackage{lipsum}
\usepackage[spanish]{babel}

\renewcommand\epigraphflush{flushright}
\renewcommand\epigraphsize{\normalsize}
\setlength\epigraphwidth{0.7\textwidth}

\definecolor{titlepagecolor}{cmyk}{1,.60,0,.40}

\DeclareFixedFont{\titlefont}{T1}{ppl}{b}{it}{0.5in}

\makeatletter                       
\def\printauthor{%                  
    {\large \@author}}              
\makeatother
\author{%
    Autor 1  \\
    Autor 2  \\
    Autor 3  \\
    Autor 4  \\
    \texttt{email1@example.com}
    }

% The following code is borrowed from: https://tex.stackexchange.com/a/86310/10898

\newcommand\titlepagedecoration{%
\begin{tikzpicture}[remember picture,overlay,shorten >= -10pt]

\node[inner sep=0pt] (UAF) at (2,0.5) {\includegraphics[width=.25\textwidth]{UAF}};
\node[inner sep=0pt] (UAF) at (-15,20) {\includegraphics[width=.25\textwidth]{uaz}};

\coordinate (aux1) at ([yshift=-15pt]current page.north east);
\coordinate (aux2) at ([yshift=-410pt]current page.north east);
\coordinate (aux3) at ([xshift=-4.5cm]current page.north east);
\coordinate (aux4) at ([yshift=-150pt]current page.north east);

\begin{scope}[titlepagecolor!40,line width=12pt,rounded corners=12pt]
\draw
  (aux1) -- coordinate (a)
  ++(225:5) --
  ++(-45:5.1) coordinate (b);
\draw[shorten <= -10pt]
  (aux3) --
  (a) --
  (aux1);
\draw[opacity=0.6,titlepagecolor,shorten <= -10pt]
  (b) --
  ++(225:2.2) --
  ++(-45:2.2);
\end{scope}
\draw[titlepagecolor,line width=8pt,rounded corners=8pt,shorten <= -10pt]
  (aux4) --
  ++(225:0.8) --
  ++(-45:0.8);
\begin{scope}[titlepagecolor!70,line width=6pt,rounded corners=8pt]
\draw[shorten <= -10pt]
  (aux2) --
  ++(225:3) coordinate[pos=0.45] (c) --
  ++(-45:3.1);
\draw
  (aux2) --
  (c) --
  ++(135:2.5) --
  ++(45:2.5) --
  ++(-45:2.5) coordinate[pos=0.3] (d);   
\draw 
  (d) -- +(45:1);
\end{scope}
\end{tikzpicture}%
}

\begin{document}
\begin{titlepage}

\noindent
\Huge Reporte de laboratorio:\\ nombre del experimento
\epigraph{Resumen:\\ \lipsum[1] }%

\null\vfill
\vspace*{1cm}
\noindent
\hfill
\begin{minipage}{0.35\linewidth}
    \begin{flushright}
        \printauthor
    \end{flushright}
\end{minipage}
%
\begin{minipage}{0.02\linewidth}
    \rule{1pt}{125pt}
\end{minipage}
\titlepagedecoration
\end{titlepage}

\newpage

\section{Introducción / Marco teórico}

\subsection{Subtema}
El subtema trata sobre lo que significa ser una ecuación numerada como la siguiente:
\begin{equation}\label{sch}
       \frac{\hbar^2}{2m}\nabla^2\Psi + V(\mathbf{r})\Psi = -i\hbar \frac{\partial\Psi}{\partial t}
\end{equation}
pero también podemos tener ecuaciones no numeradas como
\begin{equation*}
    ds^2 = c^2 dt^2 \left( \frac{d\sigma^2}{1-k\sigma^2} 
    + \sigma^2\left[ d\theta^2 
    + \sin^2\theta d\phi^2 \right] \right)
\end{equation*}
y podemos hacer referencia a cualquiera de ellas usando su etiqueta (label). Por ejemplo, 
refiriéndose a la ec. (\ref{sch}).

También podemos tener ecuaciones agrupadas y alineadas
\begin{equation}
\begin{aligned} 
2x - 5y &=  8 \\ 
3x + 9y &=  -12.
\end{aligned}
\end{equation}

o ecuaciones agrupadas por un corchete
\begin{equation}
\begin{cases}
\begin{aligned} 
2x - 5y &=  8 \\ 
3x + 9y &=  -12.
\end{aligned}
\end{cases}
\end{equation}

\subsubsection{Subsubtema }

\section{Procedimiento experimental}

\subsection{Material}

\begin{enumerate}
    \item Objeto 1
    \begin{itemize}
        \item Parte a
        \item Parte b
    \end{itemize}
    \item Objeto 2
    \item Objeto 3
\end{enumerate}

\subsection{Experimento 1}

\subsection{Experimento 2}

\section{Resultados}

\subsection{Resultados experimento 1}

\subsection{Resultados experimento 2}

\section{Análisis y discusión}

\subsection{Cuestionario}

\section{Conclusiones}

\section{Bibliografia}


\end{document}